% ============================================================
% GreenWave Networks - INFO 3180 Project
% Topic 12: Consumer Trust, Privacy, and the Ethics of AI-Native 6G Services
% AAAI 2026 Format
% ============================================================

\title{Consumer Trust, Privacy, and the Ethics of AI-Native 6G Services}
\author{Manraj Rai \and Sepehr Sheikholeslami \and Carline Gu}

\begin{document}

\maketitle

%%% -------------------------------------------------------
%%% SECTION: Abstract
%%% AUTHOR : Entire team
%%% -------------------------------------------------------
\begin{abstract}
% TODO (Entire team, finalize last): ~150 words.
% Summarize: problem (privacy/trust/ethics tension in AI-native 6G),
% scope (survey of X papers), and key contribution (proposed direction).
\end{abstract}
%%% END SECTION: Abstract


%%% -------------------------------------------------------
%%% SECTION: Introduction
%%% AUTHOR : Manraj Rai (@manrajrai50)
%%% -------------------------------------------------------
\section{Introduction}

Introducing wireless communications in the network world showed us the evolution of
how we came from the demands of modern digital applications which have increased. The
generation past us of wireless technology introduced significant improvements, from mobile
broadband in 4G networks to massive connectivity and low-latency services in 5G networks.
However, if we talk about the future now, there is growth being done with transportation, and
with health systems requiring and needing new sets of networks which are capable of managing
the extreme levels of connectivity, computation, and automation.

The 6G of network evolution is about to spiral rapidly as time goes by. The use of
artificial intelligence will be even more advanced because 6G aims to integrate AI
as a fundamental component of network operation. The overall concept is AI-native
networking, which enables the network to optimize resources, analyze conditions, and make
intelligent decisions with minimal human intervention. The research will show the
ideas of combining machine learning, distributed intelligence, and advanced communication
technologies to create highly adaptive wireless environments for the future.

With all of these advantages come challenges that also need to be addressed.
Challenges can range from privacy, security, and consumer trust to concerns around
data ownership, transparency, and responsible AI usage. Understanding the overall
foundation of AI-native 6G networks is important for seeing what solutions we can
identify for the challenges associated with future wireless services.

%%% END SECTION: Introduction


%%% -------------------------------------------------------
%%% SECTION: Background
%%% AUTHOR : Manraj Rai (@manrajrai50)
%%% -------------------------------------------------------
\section{Background: AI-Native 6G Networks}

The development of wireless communication has progressed through different generations,
each introducing new capabilities and services. While 5G introduced enhanced mobile
broadband and ultra-reliable low-latency communication, 6G aims to go further by
providing intelligent and autonomous operation, higher data rates, ultra-low latency,
massive device connectivity, and integration between communication and sensing.

AI-native 6G networks rely on several AI technologies to improve performance and
automation. Machine learning helps analyze network data, predict traffic, optimize
resources, and enhance decision-making. Deep learning supports complex wireless tasks
such as signal processing and channel estimation. Reinforcement learning allows
networks to learn from their environment and improve decisions over time. Federated
learning improves privacy by enabling devices to train AI models locally while
sharing only updates instead of raw data.

Although AI-native 6G networks offer many benefits, challenges still exist. Large-scale
data use raises privacy and user-control concerns, while AI decision-making can create
issues with transparency, fairness, and accountability. Addressing these concerns is
essential for building secure and trustworthy AI-driven 6G networks.

%%% END SECTION: Background


%%% -------------------------------------------------------
%%% SECTION: Literature Review - Privacy Risks in 6G Networks
%%% AUTHOR : Sepehr Sheikholeslami (@SEPEHRIOUS)
%%% -------------------------------------------------------
\section{Literature Review: Privacy Risks in 6G Networks}

Privacy in 6G networks has become one of the most heavily researched areas as the
network architecture moves toward AI-native operation. Nguyen et al. conducted a
comprehensive review of emerging technologies for secure communication in 6G networks,
including artificial intelligence, blockchain, quantum-safe communication, and
physical-layer security, developing a hierarchical taxonomy of security challenges
spanning physical-layer attacks, network-level threats, device vulnerabilities, and
application-specific risks \cite{nguyen2021}.

A systematic review synthesized 78 peer-reviewed studies published between 2019 and
2025 examining privacy preservation for both personal and non-personal data in 6G
networks, analyzing privacy-enhancing technologies, regulatory compliance frameworks,
and architectural patterns \cite{personaldata2026}. This connects technical privacy
mechanisms to regulatory questions such as GDPR compliance.

A more recent survey highlights how AI itself introduces new privacy risks. Emerging
applications such as LLM offloading and federated AI introduce new attack surfaces,
including prompt leakage and model inversion, while differences between regulatory
frameworks complicate cross-border data governance in multi-domain 6G systems
\cite{personaldata2026}. This shows a recurring tension in our survey's theme: the
same AI-native features marketed as privacy-preserving (e.g. federated learning) can
simultaneously open new privacy attack surfaces.

% TODO (Sepehr): Do a close reading pass of 2-3 of these papers and add
% specific methods/findings once read in full. Add corresponding
% notes/paper_summaries/*.md files for each.

%%% END SECTION: Literature Review - Privacy Risks in 6G Networks


%%% -------------------------------------------------------
%%% SECTION: Literature Review - Consumer Trust and Ethical Concerns
%%% AUTHOR : Carline Gu (@carlineGu)
%%% -------------------------------------------------------
\section{Literature Review: Consumer Trust and Ethical Concerns}

Trust has emerged as a central design requirement for AI-native 6G networks rather
than an afterthought. One survey on trust evaluation techniques for 6G networks notes
that ensuring ethical technology usage means evaluating the fairness, transparency,
and accountability of AI algorithms and data analytics to mitigate bias and unethical
conduct, and that trust evaluation should also assess the trustworthiness of
stakeholders such as network operators, service providers, and device manufacturers
\cite{trusteval2024}.

Work on trustworthy AI lifecycle management in 6G argues that a trustworthy AI
system must meet criteria including robustness, reliability, resilience,
explainability, safety, security, privacy, and fairness, and that trustworthy AI is
what ultimately drives user acceptance and confidence in autonomous network
operation \cite{reason2025}.

Industry perspectives echo this academic framing: as networks move to AI-native
architectures, trustworthiness in 6G is described as encompassing safety, security,
transparency, reliability, and ethics together, creating a paradox where AI is
necessary to manage network complexity but its black-box nature threatens the very
transparency needed to earn user trust \cite{ericsson2026}.

% TODO (Carline): Do a close reading pass of 2-3 of these papers and add
% specific methods/findings once read in full. Add corresponding
% notes/paper_summaries/*.md files for each.

%%% END SECTION: Literature Review - Consumer Trust and Ethical Concerns


%%% -------------------------------------------------------
%%% SECTION: Discussion
%%% AUTHOR : Entire team
%%% -------------------------------------------------------
\section{Discussion}

% TODO (Entire team): Synthesize the two literature review sections above.
% - Compare/contrast the privacy-focused papers vs. the trust/ethics-focused papers
% - Identify where they overlap (e.g. GDPR-related regulatory gaps appear in both)
% - Identify open gaps not addressed by current literature
% - This section should be written together once both lit reviews are finalized

%%% END SECTION: Discussion


%%% -------------------------------------------------------
%%% SECTION: Proposed Future Approach
%%% AUTHOR : Carline Gu (@carlineGu)
%%% -------------------------------------------------------
\section{Proposed Future Approach}

Current approaches to data consent in AI-native 6G networks largely mirror app-level
agreements: a user accepts a single, broad terms-of-service statement covering all data
collection and processing, without visibility into which specific AI functions are using
their data or for what purpose. As shown in the literature reviewed above, this creates
two compounding problems. First, regulatory fragmentation between frameworks such as
GDPR and PDPL complicates cross-border data governance when a single blanket consent
must satisfy multiple, sometimes conflicting, legal standards. Second, network-level
trust evaluation depends on stakeholders (operators, service providers, device
manufacturers) being individually accountable, yet a single consent agreement obscures
which stakeholder is actually responsible for a given AI-driven decision.

We propose a tiered consent model that shifts consent from the application layer to the
network function layer. Rather than a single opt-in, users would be presented with
discrete categories of AI-driven network optimization, such as energy-efficiency
management, predictive resource allocation, personalized service recommendations, and
behavioral analytics for third-party use. Users could accept or decline participation in
each category independently, and these preferences would travel with the user's profile
across network slices and providers, rather than being renegotiated at every application
boundary.

This approach offers three advantages over the status quo. First, it narrows the scope
of any single consent decision, making informed consent more meaningful rather than a
single all-or-nothing agreement. Second, it creates a natural audit trail: because each
AI function operates under its own consent category, accountability for a specific
decision can be traced to the specific function and its governing consent tier, directly
supporting the stakeholder-level trust evaluation goals raised in the trust literature.
Third, it offers a practical path to regulatory alignment, since tiers can be mapped to
the strictest applicable standard (e.g., GDPR-compliant categories) without forcing a
uniform global policy, addressing the cross-border fragmentation problem directly.

The main open challenge is technical: implementing tiered consent requires that network
functions be modular and independently identifiable at the point of data use, which is
not yet standard in current 6G architecture proposals. Realizing this approach would
likely require standardization efforts at the level of 3GPP or similar bodies, alongside
lightweight consent-signaling protocols that can operate at network speed without
introducing latency into AI-native decision loops. We view this as a promising direction
for future research rather than a fully solved problem, and a natural extension of the
privacy and trust challenges surveyed in this paper.

%%% END SECTION: Proposed Future Approach


%%% -------------------------------------------------------
%%% SECTION: Conclusion
%%% AUTHOR : Entire team
%%% -------------------------------------------------------
\section{Conclusion}

% TODO (Entire team): Summarize major findings and discuss important
% directions for future work, once all sections above are complete.

%%% END SECTION: Conclusion


%%% -------------------------------------------------------
%%% SECTION: References
%%% AUTHOR : Entire team
%%% -------------------------------------------------------
% See references.bib - AAAI/ACM style, populated as papers are read and cited above.

\end{document}
